\title{Parameter-Efficient Orthogonal Finetuning via Factorization}

\begin{abstract}
The prevalence of large foundation models has grown significantly, yet training them entirely anew remains extremely resource-intensive. Consequently, devising strategies to adapt these high-capacity models to new tasks in an efficient manner has become crucial. This work explores Orthogonal Finetuning (OFT), a principled paradigm for finetuning foundation models for downstream tasks. While OFT exhibits strong generalization, its parameter count is substantial due to the inherent high dimensionality of orthogonal matrices. To mitigate this limitation, we revisit OFT through the lens of information transfer and identify critical attributes necessary for greater parameter efficiency. Motivated by the efficient information transfer mechanisms found in the Cooley-Tukey fast Fourier transform algorithm, we introduce a compact orthogonal parameterization built upon butterfly structures. Incorporating this into the OFT framework, we devise Orthogonal Butterfly~(BOFT), a novel method for parameter-efficient finetuning. BOFT subsumes OFT as a special case, thereby providing a generalized orthogonal finetuning structure. We validate our approach with extensive empirical studies, showcasing the adaptation of large vision transformers, large language models, and text-to-image diffusion models to a diverse array of downstream vision and language tasks.
\end{abstract}

\section{Introduction}

State-of-the-art models, such as ChatGPT~\cite{brown2020language,chen2021evaluating} and Stable Diffusion~\cite{rombach2022high}, highlight the exceptional capacity of large foundation models for broad generalization. This significant advancement in model performance is accompanied by a substantial escalation in parameter scale 

(\emph{e.g.}, the GPT-3 model encompasses nearly 175B parameters). This explosion in size renders the task of training such foundation models from scratch increasingly inaccessible for many researchers. Progress in artificial intelligence research, therefore, hinges on the capability to efficiently repurpose these foundation models for new applications without the need for complete retraining. This necessity translates into the efficient adaptation of existing models for various downstream objectives. Generally, three principal categories of efficient adaptation exist: \emph{model finetuning}~\cite{hulora2022,zaken2022bitfit,guo2020parameter,raffel2020exploring,qiu2023controlling,zhang2022adaptive,chavan2023one}, where only a subset of model parameters is updated while keeping the architecture unchanged; \emph{adapter tuning}~\cite{rebuffi2017learning,houlsby2019parameter,pfeiffer2020adapterfusion,lin2020exploring,he2021towards}, which introduces trainable components into the existing model; and \emph{prompt tuning}~\cite{li2021prefix,lester2021power}, wherein additional learnable prefix tokens are prepended to the input sequence. Among these approaches, model finetuning stands out for its simplicity and effectiveness, notably without imposing extra computational overhead during inference.

Model finetuning fundamentally seeks to maintain similarity between the pretrained and finetuned models according to specific metrics, thereby safeguarding the knowledge obtained during pretraining. Commonly, contemporary finetuning strategies employ a reduced learning rate, an intuitive practice aimed at limiting divergence between the original and adapted models. Recognizing the inherent structure in weight matrices, a more systematic method attempts to retain the relational properties—specifically, the pairwise angular relationships among neurons. This consideration has inspired the development of the Orthogonal Finetuning~(OFT)~\cite{qiu2023controlling} framework, where all neurons within a layer are subjected to transformation via a single orthogonal matrix, thus ensuring that the pairwise angles among neurons are rigorously preserved throughout finetuning. While OFT exhibits strong generalizability and effective convergence in finetuning tasks such as text-to-image diffusion models~\cite{qiu2023controlling}, it encounters scalability issues: the dimensionality of orthogonal matrices leads to a substantial number of trainable parameters. OFT alleviates this by imposing a block-diagonal structure on these matrices, effectively reducing the parameter count. Nevertheless, this gain in parameter efficiency entails a trade-off—specifically, the imposed block-diagonal sparsity pattern restricts the flexibility of the transformation, as each block operates independently, potentially injecting suboptimal inductive biases (\emph{e.g.}, diminished expressiveness and inability to replicate standard linear transformations). Furthermore, determining an optimal sparsity configuration remains an unresolved challenge.

A central challenge in this context lies in constructing a dense orthogonal matrix in a manner that remains efficient with respect to parameters. At first, this appears challenging, since a dense orthogonal matrix in $d$ dimensions typically involves $\mathcal{O}(d^2)$ parameters. We introduce an alternative strategy by synthesizing a dense orthogonal matrix from a sequence of factorized sparse matrices. The foundation of this technique rests on the principle that reducing the parameter count is possible by accepting increased computational demands. Since the orthogonal matrix is realized as a product of sparse matrices, multiple multiplications are needed at every training step. To better conceptualize this matrix factorization, we reinterpret the construction of a dense orthogonal matrix as a problem of information transmission. Framed in this way, generating such a matrix through sparse multiplications corresponds to the dissemination of information across a graph with a grid topology. The information transmission viewpoint allows for the creation of diverse schemes for sparse matrix factorization that simultaneously curtail the number of learnable parameters yet possess the flexibility required to represent dense matrices.

To enhance parameter efficiency within the information transmission paradigm, our approach is informed by the butterfly patterns of the Cooley-Tukey fast Fourier transform algorithm~\cite{cooley1965algorithm}, which facilitate rapid information exchange between nodes~\cite{klasing1994broadcasting}. Specifically, the butterfly network in the Cooley-Tukey algorithm inspires an efficient form of sparse matrix factorization, referred to as butterfly factorization. Utilizing this structure, a dense matrix of size $d$ can be expressed as the product of $\mathcal{O}(\log d)$ sparse matrices, each containing only $\mathcal{O}(d)$ non-zero elements. By maintaining orthogonality in each sparse matrix, this leads to a highly efficient orthogonal parameterization requiring merely $\mathcal{O}(d\log d)$ parameters—representing a substantial savings relative to the standard approach. On the strength of this compact orthogonal parameterization, we introduce a novel, parameter-efficient finetuning strategy: Orthogonal Butterfly~(BOFT). BOFT extends OFT as a specific instance and offers a comprehensive orthogonal finetuning architecture. Notably, both the block-diagonal and butterfly constructions share an underlying feature—namely, sparsity—leveraging this property to limit the parameter count in orthogonal matrices. Comparing our method with the low-rank framework seen in LoRA~\cite{hulora2022} is instructive; both low-rank and sparse matrices are principal structured matrix types~\cite{candes2011robust} that promote efficient use of parameters.

In contrast to the block-diagonal architecture employed by OFT, which manages the trade-off between expressivity and regularity, BOFT adopts the butterfly structure to facilitate a more continuous transition between full orthogonal matrices (high expressivity) and the identity matrix (strong regularity). This approach permits the identification of a narrower set of candidate matrices within the orthogonal group tailored to downstream applications. Given that the butterfly structure underpins numerous efficient linear operators, including the discrete Fourier and discrete cosine transforms, we contend that employing such structured parameterization induces an inherent inductive bias. This bias is advantageous for generalization and reduces the risk of overfitting. Our intuition rests on the fact that the butterfly structure can naturally recover a variety of classical linear transforms—a property absent in the block-diagonal configurations used by OFT. The main contributions of our work are summarized as follows:

\begin{itemize}[leftmargin=*,nosep]
    \item We address parameter efficiency in the context of orthogonal finetuning by introducing a new information transmission framework, which reconceptualizes the construction of dense, parameter-efficient orthogonal matrices as an information propagation challenge on a grid-structured graph.
    \item Drawing on the butterfly architecture from the Cooley-Tukey algorithm, we introduce Orthogonal Butterfly—a parameter-efficient method for orthogonal finetuning—within the information transmission paradigm.
    \item We present several theoretical justifications for BOFT’s ability to sharply decrease the number of trainable parameters, yet maintain strong expressivity and stability during training. Due to its underlying matrix factorization, BOFT further enables a notable weight interpolation property (see Figure~\ref{fig:boft_interpolation}).
    \item To our knowledge, this is the inaugural exploration of orthogonal finetuning in settings beyond controllable text-to-image generation~\cite{qiu2023controlling}, thereby showcasing its promise as a general-purpose model adaptation technique. To validate this, we apply BOFT across a diverse array of adaptation tasks, including both computer vision and natural language processing, and consistently observe that it outperforms contemporary state-of-the-art approaches in terms of parameter-efficiency and generalization.
\end{itemize}

\section{Related Work}

\textbf{Parameter-efficient finetuning (PEFT)}. The trend toward ever-larger and more capable foundation models (\emph{e.g.}, \cite{brown2020language,radford2021learning,rombach2022high,kirillov2023segment}) has led to a surge of research on how to adapt such models to new tasks with minimal additional parameters~\cite{treviso2023efficient,ding2023parameter,lialin2023scaling}. Of the many approaches categorized as PEFT~\cite{houlsby2019parameter,aghajanyan2020intrinsic,hulora2022,edalati2022krona,wang2022adamix,gheini2021cross,zaken2022bitfit,guo2020parameter,sung2021training,ansell2022composable,lester2021power,li2021prefix,vu2022spot,he2021towards,mao2021unipelt,karimi2021compacter,liu2022few,sung2022lst,chen2022parameter,jia2022visual,chen2022adaptformer,zhang2022neural,jie2023fact,lian2022scaling,luo2023towards}, those based on reparameterization~\cite{aghajanyan2020intrinsic,hulora2022,edalati2022krona} are most closely related to our work. LoRA~\cite{hulora2022} introduces updates to pretrained weights using a pair of low-rank matrices whose product is added to the original matrix, leading to strong results on NLP benchmarks. However, LoRA uses a fixed rank for all additional matrices, motivating methods such as \cite{zhang2022adaptive,valipour2022dylora,zhang2023increlora} that allow for adaptive rank allocation per layer to better fit parameter constraints. Because of its straightforward design, this low-rank modification has seen widespread adoption \cite{dettmers2023qlora,zi2023delta,chavan2023one}. Drawing from insights on how hyperspherical energy informs generalization properties~\cite{liu2018learning,liu2021orthogonal}, \cite{qiu2023controlling} developed orthogonal finetuning (OFT)—a different yet effective scheme for tuning text-to-image diffusion models. Here, OFT trains an orthogonal transformation that operates on neurons within a layer, consistently producing more stable optimization and stronger generalization than LoRA. Nonetheless, the number of trainable parameters in OFT typically exceeds that of LoRA, making improvements in OFT’s parameter efficiency desirable. Furthermore, the applicability of OFT to broader adaptation tasks beyond text-to-image diffusion~\cite{qiu2023controlling} has not been systematically explored. BOFT advances this direction by employing butterfly factorization to compress OFT’s parameterization. This approach enables us to extend and validate the benefits of orthogonal finetuning across general adaptation problems.

\textbf{Butterfly structures}. The radix-2 Cooley-Tukey algorithm is known for recursively breaking down the $N$-point discrete Fourier transform into smaller $\frac{N}{2}$-point subproblems, yielding a characteristic butterfly pattern expressible as products of sparse matrices (also referred to as butterfly matrices). Such matrices have found utility in various settings, such as serving as orthogonal matrix parameterizations to eliminate the need for pivoting in Gaussian elimination and boosting computational efficiency~\cite{parker1995random}, stabilizing the training process for recurrent neural architectures~\cite{jing2017tunable}, and approximating kernels~\cite{munkhoeva2018quadrature}. Techniques introduced in \cite{dao2019learning,dao2019kaleidoscope} learn fast linear mappings based on butterfly matrix structures. Butterfly matrices have further enabled efficient neural network training~\cite{chen2022pixelated,dao2022monarch}. Their computational advantages also extend to accelerating matrix-vector multiplication~\cite{michielssen1996multilevel,o2010algorithm}, approximating matrices in data-sparse regimes~\cite{li2015butterfly}, and supporting network broadcast or data transmission~\cite{klasing1994broadcasting,li2003linear,rajasingh2016transmission}. Distinct from previous efforts, our objective is to leverage butterfly structures to enhance OFT’s parameter efficiency.

We begin by outlining the foundational concepts of OFT. In order to adapt a pretrained weight matrix, OFT reparameterizes the updated weight matrix as the product of a learnable orthogonal matrix and the original, fixed weight matrix. This approach contrasts with LoRA, which employs an additive low-rank update to modify the weights, whereas OFT utilizes a multiplicative update via an orthogonal matrix. Parameter-efficiency in LoRA is achieved through its low-rank design, while the initial formulation of OFT incorporates a block-diagonal orthogonal matrix structure~\cite{qiu2023controlling}, with smaller block sizes enhancing parameter efficiency. Figure~\ref{fig:comp_lora} provides a visual comparison. The key rationale for applying an orthogonal transformation when finetuning is to maintain the relative angles between neuron vectors~\cite{liu2018learning,liu2021orthogonal,qiu2023controlling}, thereby preserving much of the semantic information acquired during pretraining.

More formally, OFT learns an orthogonal matrix $\bm{R}\in\mathbb{R}^{d\times d}$ to augment a pretrained linear layer with weights $\bm{W}^0\in\mathbb{R}^{d\times n}$, modifying the layer’s forward computation from $\bm{z}=(\bm{W}^0)^\top\bm{x}$ to $\bm{z}=(\bm{R}\bm{W}^0)^\top\bm{x}$, where $\bm{x}\in\mathbb{R}^{d}$ and $\bm{z}\in\mathbb{R}^{n}$ are the respective input and output. For zero initialization, $\bm{R}$ is set as the identity matrix initially. To guarantee $\bm{R}$ remains orthogonal during finetuning, the Cayley parameterization is adopted, following~\cite{qiu2023controlling,liu2021orthogonal}: specifically, $\bm{R}=(\bm{I}+\bm{Q})(\bm{I}-\bm{Q})^{-1}$ where $\bm{Q}$ is a skew-symmetric matrix such that $\bm{Q}=-\bm{Q}^\top$. For efficient parameterization, the orthogonal matrix $\bm{R}$ is structured as a block-diagonal matrix $\text{diag}(\bm{R}_1, \bm{R}_2, \cdots, \bm{R}_r)$, in which each $\bm{R}_i\in\mathbb{R}^{b\times b}$ is an orthogonal matrix and $br=d$. This block-diagonal configuration reduces the parameter count, but at the expense of introducing the constraint that the vector components of each neuron (i.e., the columns of $\bm{W}^0$) are arbitrarily partitioned into $r$ groups and independently transformed by separate orthogonal matrices. Although this design yields empirical success, there is little theoretical justification for partitioning the neuron dimensions solely by their indices, making a dense orthogonal matrix more appealing. This observation prompts a critical question: \emph{Is it feasible to construct a dense orthogonal matrix while retaining parameter efficiency?}

To tackle this issue, we suggest constructing a dense orthogonal matrix as the product of several sparse orthogonal matrices. To present a coherent and intuitive framework for analyzing the sparsity structures involved in orthogonal matrix factorization, we conceptualize the task of synthesizing a dense orthogonal matrix in OFT through the lens of information transmission. Concretely, when forming a dense matrix $\bm{R}\in\mathbb{R}^{d\times d}$ by multiplying $m$ square matrices, $\thickmuskip=2mu \medmuskip=2mu \bm{R}=\bm{B}_m\bm{B}_{m-1}\cdots\bm{B}_1$, we can interpret this process as disseminating information across a $d\times (m+1)$ node grid, as depicted in Figure~\ref{fig:info_trans}. This perspective is motivated by noting that a $d$-dimensional dense square matrix encodes all-to-all connections from $d$ source nodes to $d$ target nodes. If for matrix $\bm{R}$ the entry $R_{ij}$ equals zero, it signifies an absence of information transfer from the $j$-th input node to the $i$-th output node; conversely, a non-zero entry indicates feasible information transmission. Therefore, representing $\bm{R}$ by the sequential multiplication $\bm{B}_m\bm{B}_{m-1}\cdots\bm{B}_1$ amounts to staging successive rounds of information exchange as governed by the graph structures determined by each $\bm{B}_i$. Information propagation initiates according to the topology specified by $\bm{B}_1$, proceeding sequentially through to $\bm{B}_n$. As an illustrative case in Figure~\ref{fig:info_trans}, we examine the decomposition $\bm{R}=\bm{B}_5\bm{B}_4\bm{B}_3\bm{B}_2\bm{B}_1$, with its associated sparsity patterns and induced connectivity visualized. The corresponding graph results from an explicit unrolling of the matrix multiplications. Within this induced structure, each $\bm{B}_i$ specifies the connections from nodes at level $i$ to those at level $i+1$: the $(j_1,j_2)$ entry in $\bm{B}_i$ determines whether there exists a directed edge from the $j_2$-th node at level $i$ to the $j_1$-th node at level $i+1$ (zero indicating no such connection). For $\bm{B}_5\bm{B}_4\bm{B}_3\bm{B}_2\bm{B}_1$ to yield a dense outcome, each node at the initial level must be able to reach every node at the sixth level. Examining the reduced case $\bm{R}=\bm{B}_4\bm{B}_3\bm{B}_2\bm{B}_1$—connecting first-level sources to fifth-level receivers—we observe, for instance, that information originating from node $1$ cannot reach node $3$, reflecting a zero entry at position $(3,1)$ in the sparsity pattern of $\bm{B}_4\bm{B}_3\bm{B}_2\bm{B}_1$. This relationship between the graph representation and the sparsity structure persists when considering further truncated products such as $\bm{R}=\bm{B}_3\bm{B}_2\bm{B}_1$ or $\bm{R}=\bm{B}_2\bm{B}_1$. Broadly, to ensure $\bm{R}\in\mathbb{R}^{d\times d}$ is fully dense, the set of factors $\bm{B}_m,\ldots,\bm{B}_1$ must define a directed edge configuration on the $d\times (m+1)$ grid so that each directed link connects a pair of nodes between adjacent columns (levels), ensuring that every source at the first level is able to transmit to any destination node at the $(m+1)$-th level.

Adopting the information transmission perspective offers valuable insight into the sparsity structure of the factorization matrices used in OFT. In Figure~\ref{fig:blk_diag}, the block-diagonal arrangement of $\bm{R}$ in the standard OFT framework is illustrated. While this configuration decreases the count of trainable parameters, it is inherently incapable of producing a dense matrix $\bm{R}$. Our objective is to synthesize a dense orthogonal matrix by combining $m$ sparse orthogonal matrices, aiming to do so with the minimum number of effective trainable parameters. From the information transmission standpoint, our main requirements can be summarized as follows: \emph{(i)} \textbf{dense connectivity}: every node in the initial layer must be connected via at least one path to every node in the terminal layer; and \emph{(ii)} \textbf{minimal free edges}: the aggregate number of edges should be minimized, subject to orthogonality. Orthogonality places a significant constraint on the permissible edges between adjacent layers. Specifically, to ensure that each matrix $\bm{B}_i$ maintains full rank—a prerequisite for orthogonality—it is essential to establish $d$ edges, enabling a one-to-one correspondence between nodes at the $i$-th level and those at the $(i-1)$-th level. This implies a lower bound of $d$ edges between consecutive layers (\emph{e.g.}, 4 as shown in Figure~\ref{fig:info_trans}). Such $d$ edges are obligatory for orthogonality and should be excluded when tallying the number of edges, since these elements remain fixed and untrainable (\emph{e.g.}, in a $d\times d$ orthogonal matrix with $d$ non-zero elements, those values are confined to $\pm 1$). Due to orthogonality reducing the parameter space compared to dense matrices, it introduces interdependencies among the edges. For instance, in a $2\times 2$ orthogonal matrix, setting the entry at $(1,1)$ to zero necessitates a corresponding zero at $(2,2)$ (\emph{i.e.}, the absence of one edge compels the absence of another). Thus, orthogonality may require the addition or removal of certain edges for each configuration. Examining the pattern of non-zero entries in the constructed orthogonal matrix, the information transmission model is especially informative in OFT, since our focus is solely on which entries of $\bm{R}$ are trainable and non-zero, rather than their actual values.

A straightforward approach to establish dense connections between two layers necessitates $\mathcal{O}(d^2)$ edges (that is, forming a single dense orthogonal matrix), which comprises $d^2-d$ connections (for instance, when $d=4$, there are 12 edges involved). Figure~\ref{fig:info_trans} presents an illustrative example of a practical matrix factorization, which requires only $10$ edges—fewer than needed for a single dense orthogonal matrix. This setup allows for analysis of the parameter efficiency of OFT from the perspective of graph structures, making it convenient to conceptualize valid matrix factorizations within this paradigm. Inspiration is drawn from a particular network topology in the Cooley-Tukey algorithm, specifically the butterfly graphs~\cite{cooley1965algorithm}, which provide a highly efficient way to fully connect $d$ source nodes to $d$ receiver nodes, utilizing just $\mathcal{O}(d\log d)$ connections. As demonstrated in Figure~\ref{fig:info_trans}, the depicted topology accomplishes dense connectivity with $10$ edges, yet a butterfly network achieves this with only $8$ connections. In what follows, we detail how the butterfly structure can enhance the efficiency in the number of parameters.

\section{Orthogonal Parameterization by Butterfly Factorization}

Initially developed in the context of the Cooley-Tukey algorithm for efficient fast Fourier transform computation, the butterfly structure facilitates scenarios where localized modifications in the frequency domain manifest as widespread changes in the spatial domain. This dynamic is analogous to our context of information transmission: every node from the input layer is able to propagate information to every output node in the final layer. The butterfly pattern has also been adopted as an efficient design for computer network topologies~\cite{solihin2015fundamentals,li2011linear}, facilitating effective communication. For any $k \geq 2$ where $k$ is a power of $2$, we introduce the \emph{butterfly factor} $\bm{B}^F(k)$ defined as

\begin{equation}\label{eq:but_factor}
\begin{aligned}
    \bm{B}^F(k)=\begin{bmatrix}
    \text{diag}(\bm{d}_1) & \text{diag}(\bm{d}_2) \\
    \text{diag}(\bm{d}_3) & \text{diag}(\bm{d}_4)
    \end{bmatrix}\in\mathbb{R}^{k\times k},
\end{aligned}
\end{equation}

where each vector $\bm{d}_i\in\mathbb{R}^{\frac{k}{2}}$ for all $i$. Subsequently, for $d=2^N$, the $d$-dimensional \emph{butterfly matrix} $\bm{B}(d)\in\mathbb{R}^{d\times d}$ can be described recursively as follows:

\begin{equation}
\begin{aligned}
    \!\!\bm{B}(d)=\tilde{\bm{B}}(d,d)\!\cdot\!\begin{bmatrix}
    \bm{B}_1(\frac{d}{2})\!\!\!\!\! & \bm{0} \\
    \bm{0}\!\!\!\!\! &\bm{B}_2(\frac{d}{2})
    \end{bmatrix}= \tilde{\bm{B}}(d,d)\tilde{\bm{B}}(d,\frac{d}{2})\cdots\tilde{\bm{B}}(d,2),
\end{aligned}
\end{equation}

where $\bm{B}_1(\frac{d}{2})$ and $\bm{B}_2(\frac{d}{2})$ represent two butterfly matrices of dimension $\frac{d}{2}$. We introduce the concept of the \emph{butterfly component} as $\thickmuskip=2mu \medmuskip=2mu \tilde{\bm{B}}(d,k)=\text{diag}(\bm{B}^F_1(k),\cdots,\bm{B}^F_{d/k}(k))$, which is constructed as a block-diagonal matrix of size $d\times d$ with each block having size $k$, where each $\bm{B}^F_i(k)$ is a butterfly factor as outlined in Equation~\ref{eq:but_factor}. We can now utilize butterfly matrices for the orthogonal parameterization. To do so, it is sufficient to ensure that every butterfly component $\tilde{\bm{B}}(d,k)$ in $\bm{B}(d)$ maintains orthogonality. Focusing first on the block-diagonal matrix $\tilde{\bm{B}}(d,2)$ with block size $2$, orthogonality can be straightforwardly achieved by applying the Cayley transform (or $2$-dimensional rotation) to parameterize each individual block, following the methodology of \cite{liu2021orthogonal,qiu2023controlling}. The sparsity structure of each butterfly component is simply a permuted version of that of $\tilde{\bm{B}}(d,2)$, enabling all butterfly components to be parametrized as orthogonal matrices in a similar fashion. This framework allows us to efficiently represent orthogonal matrices using a composition of $2\times 2$ orthogonal blocks. Extending the butterfly matrices according to \cite{chen2022pixelated}, we define the block butterfly component $\tilde{\bm{B}}^b(d,k)$, where every entry in $\bm{d}_i$ is now a $b\times b$ block. To preserve orthogonality for the block butterfly component $\tilde{\bm{B}}^b(d,2)$, each $2b\times 2b$ block matrix is independently parameterized as an orthogonal matrix. For butterfly components $\tilde{\bm{B}}^b(d,k)$ with $k>2$, their non-zero block structures can be seen as permuted forms of those in $\tilde{\bm{B}}^b(d,2)$, thus allowing a similar orthogonal parameterization. Assembling these elements, the forward propagation step in BOFT is given by

\begin{equation}
    \begin{aligned}\nonumber
    \bm{z}=\big(\bm{R}(m,b)\cdot\bm{W}^0\big)^\top\bm{x},~~~~\text{s.t.}~\bigg{\{}\bm{R}(m,b)=\prod_{i=1}^m \tilde{\bm{B}}^b_{(d,i)}~~\&~~\underbrace{\big(\tilde{\bm{B}}^b_{(d,j)}\big)^\top\tilde{\bm{B}}^b_{(d,j)}=\tilde{\bm{B}}^b_{(d,j)}\big(\tilde{\bm{B}}^b_{(d,j)}\big)^\top\!=\bm{I}_d}_{\forall j\in [1, m]}\bigg{\}},
    \end{aligned}
\end{equation}

Here, for clarity, $\tilde{\bm{B}}^b(d,2^{m - i+1})$ is denoted by $\tilde{\bm{B}}^b_{(d,i)}$, and $\bm{I}_d$ denotes the $d$-dimensional identity matrix. The orthogonal transformation $\bm{R}(m,b)\in\mathbb{R}^{d\times d}$ is built as a sequential product of multiple butterfly matrices, each being orthogonal. For notational convenience, BOFT employing $\bm{R}(m,\frac{b}{2})$ is referred to as BOFT($m$,$b$), provided $b\geq 2$. When $\thickmuskip=2mu \medmuskip=2mu m=1$, BOFT($1$,$b$) corresponds to the block-diagonal OFT~\cite{qiu2023controlling} using blocks of size $b$. Additionally, setting $b=d$ yields BOFT($1$,$d$), which matches the original OFT~\cite{qiu2023controlling} where the orthogonal matrix is unconstrained and dense. If we set $m=\log \frac{2d}{b}$, BOFT employs the block butterfly matrix $\bm{B}^{\frac{b}{2}}(d)$ as the orthogonal transformation $\bm{R}$, thereby producing a dense orthogonal matrix. Generally, BOFT($m$,$b$) utilizes $\frac{1}{2}(b-1)dm$ effective learnable parameters to adapt a linear transformation of dimensions $d\times n$. When the butterfly structure is used, i.e., $m=\log d,b=2$, the parameter count in BOFT scales as $\mathcal{O}(d\log d)$. By comparison, standard OFT based on a fully dense orthogonal matrix involves $\mathcal{O}(d^2)$ parameters, while a block-diagonal OFT with $r$ blocks incurs $\mathcal{O}(bd)$ parameters. Thus, conventional OFT must set $b=d$ to obtain a dense orthogonal matrix, but BOFT can achieve a dense orthogonal transformation regardless of the block size $b$.

\textbf{Identity Initialization in BOFT}. In most finetuning approaches, initialization aligns with the exact pretrained weights, ensuring that further training does not significantly diverge from the original model. For instance, LoRA initializes the low-rank weight perturbation at zero. In the context of BOFT, every butterfly block is initialized as an identity matrix (meaning the skew-symmetric matrices in the Cayley transform start from zero).

\textbf{Multiplicative Dropout in BOFT}. To enhance regularization and combat overfitting, LoRA~\cite{hulora2022} applies a Dropout mechanism to the low-rank parameter updates. Traditional Dropout~\cite{srivastava2014dropout} is naturally compatible with LoRA's additive updates, but is not directly applicable to BOFT's multiplicative updates. To resolve this, we introduce a multiplicative Dropout specifically designed for BOFT: Since $\bm{R}(m,b)$ consists of $m$ orthogonal butterfly modules, which may be permuted into $2b\times 2b$ block-diagonal forms, our method randomly selects $p_1$ proportion of butterfly modules and within each, a $p_2$ fraction of diagonal sub-blocks, and sets the chosen blocks to identity matrices.

\section{Intriguing Insights and Discussions}

\textbf{Expressivity of BOFT}. The butterfly architecture combined with permutations is capable of exactly recovering various canonical fast linear transforms~\cite{dao2019learning,dao2019kaleidoscope} (\emph{e.g.}, the fast Fourier transform, Hadamard transform). Nevertheless, it remains an open question how accurately our orthogonal butterfly matrix can represent arbitrary orthogonal matrices. To investigate this, we perform an experiment aiming to approximate a randomly sampled dense orthogonal matrix~\cite{anderson1987generation} of dimensions $512\times 512$. As shown in Figure~\ref{fig:para_eff}, results are averaged across 10 independent random initializations. On the vertical axis, we plot the approximation error, while the horizontal axis corresponds to the number of effective trainable parameters. For each block size, BOFT is depicted by a distinct color, with the leftmost point marking the performance of BOFT($1$,$b$) (i.e., the baseline block-diagonal OFT with block size $b$). Across configurations, BOFT generally demonstrates superior parameter efficiency in comparison to OFT. For instance, BOFT($9$,$2$) surpasses BOFT($1$,$16$) in expressivity despite using considerably fewer parameters. Moreover, employing a smaller $b$ with a higher $m$ tends to improve parameter efficiency; as an illustration, BOFT($6$,$4$) uses substantially fewer parameters yet achieves comparable approximation errors to BOFT($2$,$16$). In summary, the butterfly matrix forms a more structured subset within the orthogonal group compared to the block-diagonal configuration, thus ensuring that BOFT is provably more expressive than OFT for equivalent block sizes.

\begin{theorem}[Expressivity of BOFT]\label{thm:exp_boft}
    BOFT possesses greater expressivity than OFT under equal block size. Any orthogonal matrix of size $d$ can be closely approximated by composing butterfly matrices of the form $\bm{B}_{d-1,1}(d)\bm{B}^\top_{d-1,2}(d)\cdots\bm{B}_{1,1}(d)\bm{B}^\top_{1,2}(d)$, where each $\bm{B}_{i,j}(d)$ denotes a butterfly matrix for all relevant $i$ and $j$.
\end{theorem}

Building on Theorem~\ref{thm:exp_boft}, one can straightforwardly extend BOFT by representing the final orthogonal matrix as $\thickmuskip=2mu \medmuskip=2mu \bm{R}^G(m_1,b_1,m_2,b_2,l)=\bm{R}_{l,1}(m_1,b_1)\bm{R}_{l,2}^T(m_2,b_2)\cdots\bm{R}_{1,1}(m_1,b_1)\bm{R}_{1,2}^T(m_2,b_2)$, where $\bm{R}_{i,j}^T(m,b)$ refers to the orthogonal matrices parameterized in BOFT. Setting $m_1=m_2=\log d$, $b_1=b_2=2$, and $l=d-1$ allows $\bm{R}^G(m_1,b_1,m_2,b_2,l)$ to encompass the entire orthogonal group, a construction known as the kaleidoscope hierarchy~\cite{dao2019kaleidoscope}. Still, enhanced expressivity does not inherently guarantee superior downstream performance in finetuning tasks—indeed, full finetuning, despite its complete expressivity, is often outperformed by more regularized approaches. Striking a balance between expressivity and structural regularity is crucial for achieving robust generalization during model finetuning. By incorporating structured priors, BOFT expands the accessible finetuning parameter space, thereby facilitating an improved balance between expressivity and regularity.

\textbf{Spectral properties}. Orthogonal finetuning typically results in superior spectral characteristics compared to LoRA, due to its exact preservation of the spectral norm of the pretrained weight matrix $\bm{W}^0$. This can be demonstrated using singular value decomposition, where $\bm{W}^0=\bm{U}\bm{\Sigma}\bm{V}^\top$ and $\bm{U},\bm{V}$ are orthogonal matrices while $\bm{\Sigma}$ denotes the diagonal matrix of singular values. In both OFT and BOFT, an orthogonal matrix $\bm{R}$ is applied on the left, resulting in finetuned parameters $\bm{R}\bm{U}\bm{\Sigma}\bm{V}^\top$; this operation does not modify the largest singular value (that is, the spectral norm of $\bm{W}^0$). Maintaining this property has been observed to significantly enhance both training stability and model generalization~\cite{miyato2018spectral,yoshida2017spectral}. Additional mathematical insights regarding this property are provided in Appendix~\ref{app:sn_property}.

\textbf{Orthogonal finetuning as learning bilinear similarity}. The BOFT approach can be interpreted as learning the bilinear similarity $\bm{w}^0_i\bm{R}\bm{x}$, where $\bm{w}^0_i$ denotes the $i$-th neuron (specifically, the $i$-th column of $\bm{W}^0$). In this view, BOFT involves learning the matrix $\bm{R}$ governing bilinear similarity, under the constraint that $\bm{R}$ remains orthogonal; this connects closely to frameworks such as distance metric learning~\cite{xing2002distance} and the theory of bilinear forms~\cite{roman2005advanced}.

\textbf{Inductive bias and generalization in BOFT}. The $\bm{R}(m,b)$ matrix in BOFT is usually confined to a structured subset within the orthogonal group, effectively restricting the hypothesis space and inducing an inductive bias. We propose that the specific structure enforced by butterfly factorization imparts an inductive bias that is advantageous for generalization, since it reflects patterns shared by various well-established linear transforms~\cite{dao2019kaleidoscope}, such as the discrete Fourier transform, the discrete sine/cosine transforms, and the Hadamard transform. Furthermore, the sparse matrix factorization embedded in BOFT might introduce additional implicit forms of inductive bias~\cite{gunasekar2017implicit,li2018algorithmic,liu2019neural}.

\textbf{Comparison to butterfly-based sparse training}. Numerous studies~\cite{dao2019learning,dao2019kaleidoscope,chen2022pixelated,dao2022monarch} have investigated sparse training utilizing butterfly parameterizations. These approaches generally involve directly reparameterizing the weight matrices with butterfly structures and training the networks from initialization. In \cite{dao2022monarch}, the authors examine the adaptation of pretrained weights by initially projecting them onto modified butterfly matrices, followed by optimizing these projected parts for specific downstream applications. In contrast, BOFT introduces an alternative finetuning technique, employing layer-shared weight matrices to transform weights, which differs fundamentally from the approaches taken by prior work.

\input{rebuttal-results}

\section{Concluding Remarks and Limitations}

In this work, we introduce Orthogonal Butterfly, a broadly applicable parameter-efficient finetuning approach for foundation models that leverages butterfly structures. The central idea behind improving parameter efficiency is to express a dense orthogonal matrix as a product of multiple sparse orthogonal matrices. To facilitate the discovery of suitable matrix decompositions, we design a graphical information transmission framework. Within this conceptual structure, we identify that the butterfly architecture is particularly well-suited for our goals of sparse orthogonal matrix factorization. We present empirical results demonstrating the utility of BOFT for finetuning across large-scale language, vision, and text-to-image generative models, showing that BOFT consistently outperforms other generic finetuning methods.

Nevertheless, there are certain limitations to BOFT. The requirement to represent the final orthogonal matrix as a product of several orthogonal matrices incurs a modest increase in training runtime compared to OFT. Addressing this overhead in BOFT's training remains an unresolved challenge. Fortunately, after finetuning, the orthogonal matrices learned by BOFT can be merged into the original model, resulting in no extra latency at inference time. Another open question concerns whether butterfly networks offer the optimal structure for efficient information transmission. Our information transmission framework opens the door for further insights from the field of computer networking, where the efficiency of network topologies for data transmission is rigorously explored. This connection suggests that alternative and possibly more efficient structures might exist for constructing dense orthogonal matrices.

\end{document}